% !PRINTED 778
\weberpage{778}
% running head "II. Nachtrag." at the top of this page is not transcribed as body text.

Here, however, we have no need to make use of the special expressions for these
characters. We rely instead on the following theorem, which holds for all
\gesperrt{Abel}ian groups:

1. \gesperrt{If} $A$ \gesperrt{is an element of degree} $f$ \gesperrt{of an}
\gesperrt{Abel}\gesperrt{ian group of degree} $\mu$, \gesperrt{and if} $\mu = ef$,
\gesperrt{then all the} $\chi_i(A)$ \gesperrt{are} $f$\gesperrt{th roots of unity,
and among them every} $f$\gesperrt{th root of unity occurs exactly} $e$
\gesperrt{times.}

The theorem is contained implicitly in theorem 7., \S.\,12 of this volume. For if we
apply that theorem to the group
\[
T = 1, A, A^2, \ldots, A^{f-1}
\]
whose index with respect to $\mathfrak{N}$ is equal to $e$, it follows that there is
a group $\Xi$ of $e$ characters $\xi$ for which the
conditions
\[
\xi(A) = 1
\]
are satisfied.

If, then, one decomposes the group $X$ of the characters $\chi$ into the subsidiary
groups
\[
\Xi_1, \Xi_2, \ldots, \Xi_f,
\]
then $\chi_1(A)$ and $\chi_2(A)$ are equal to one another or different from one
another according as $\chi_1$ and $\chi_2$ occur in the same or in different ones of
these subsidiary groups. Since, moreover, because
$\chi(A)^f = \chi(A^f) = 1$ all the $\chi(A)$ are roots of unity of degree $f$, and
there are only $f$ different such roots, our theorem is thereby proved.

If $n$ is contained in the class $A$, then the degree $f$ of $A$ is the exponent to
which $n$ \gesperrt{belongs} with respect to the modulus $m$, that is, the least positive
exponent for which
\weq{5}{n^f \equiv 1 \pmod{m}}
holds. If then $n$ belongs to the exponent $f$, all the $\chi_i(n)$ are $f$th roots
of unity, and every $f$th root of unity occurs $e$ times among them.

If $a$ denotes any one of the $\varphi(m)$ residues of the numbers $n$, then in the
form
\weq{6}{\mu_x = mx + a - m}
every number of the number class (2) represented by $a$ is contained; if by $a$ we
understand the least positive residue, then we obtain the totality of the numbers of
this class when we let $x$ run through all positive integral values.
